% --- Archon physics formalization source begin ---
% archon:physics
% archon:covers PhyXMiniProblems/problem_phyx_mini_0647.lean
% archon:source-report reports/phyx_mini/problem_phyx_mini_0647.source.json

\chapter{Physics problem phyx\_mini\_0647}
\label{ch:PhyXMiniProblems_problem_phyx_mini_0647}

\paragraph{Problem source.}
\#\# Physical scenario

Figure shows the wave function of an electron.

\#\# Question

What is the probability that the electron is located between \$x = -1.0\$ \$nm\$ and \$x = 1.0\$ \$nm\$?

\#\# Answer choices

- A: A: 0.40

- B: B: 0.80

- C: C: 0.90

- D: D: 0.60

\#\# Figure caption (auxiliary; use the image as primary evidence)

The image is a graph plotting a mathematical function \textbackslash{}(\textbackslash{}psi(x)\textbackslash{}) against \textbackslash{}(x\textbackslash{}) (in nanometers). 

\#\#\# Details:

- **Axes:**
  - The horizontal axis is labeled \textbackslash{}(x (\textbackslash{}text\{nm\})\textbackslash{}).
  - The vertical axis is labeled \textbackslash{}(\textbackslash{}psi(x)\textbackslash{}).

- **Scale and Labels:**
  - The \textbackslash{}(x\textbackslash{})-axis has tick marks at \textbackslash{}(-2\textbackslash{}), \textbackslash{}(-1\textbackslash{}), \textbackslash{}(1\textbackslash{}), and \textbackslash{}(2\textbackslash{}).
  - The \textbackslash{}(\textbackslash{}psi(x)\textbackslash{})-axis is marked with \textbackslash{}(-c\textbackslash{}), \textbackslash{}(-\textbackslash{}frac\{1\}\{2\}c\textbackslash{}), \textbackslash{}(\textbackslash{}frac\{1\}\{2\}c\textbackslash{}), and \textbackslash{}(c\textbackslash{}).

- **Function Characteristics:**
  - The graph is composed of several horizontal line segments:
    - From \textbackslash{}(-2\textbackslash{}) to \textbackslash{}(-1\textbackslash{}), and from \textbackslash{}(1\textbackslash{}) to \textbackslash{}(2\textbackslash{}), the value of \textbackslash{}(\textbackslash{}psi(x)\textbackslash{}) is \textbackslash{}(-c\textbackslash{}).
    - From \textbackslash{}(-1\textbackslash{}) to \textbackslash{}(1\textbackslash{}), the value of \textbackslash{}(\textbackslash{}psi(x)\textbackslash{}) is \textbackslash{}(c\textbackslash{}).
    - At \textbackslash{}(x = -1\textbackslash{}) and \textbackslash{}(x = 1\textbackslash{}), the function has vertical lines, transitioning between \textbackslash{}(-c\textbackslash{}) and \textbackslash{}(c\textbackslash{}).

This diagram appears to illustrate a piecewise function with different constant values over distinct intervals.

\#\# Dataset metadata

Quantum Phenomena; Physical Model Grounding Reasoning, Numerical Reasoning

\paragraph{Recorded answer/context.}
B: B: 0.80

\paragraph{Figure/image path.}
/root/proposal\_for\_physic/science-mango/phyx\_mini\_run/phyx\_data/test\_image/647.png

\paragraph{Formalization target.}
create a compiling Lean file with sorry bodies at `PhyXMiniProblems/problem\_phyx\_mini\_0647.lean`. The Lean declarations must preserve the physical quantities, dimensions or dimensional roles, figure labels, governing-law hypotheses, and final relation expressed by this problem.
Use Mathlib/Physlib names found through LeanExplore where available. If a physics API is missing, introduce faithful local abstractions rather than scalar placeholder aliases.

\begin{theorem}[Physics formalization target]
\label{thm:physics:phyx_mini_0647:target}
\lean{PhyXMiniProblems.ProblemPhyXMini0647.problem_phyx_mini_0647}
\uses{def:physics:phyx-mini-0647:phyxminiproblems-problemphyxmini0647-electronpositionexperiment, def:physics:phyx-mini-0647:phyxminiproblems-problemphyxmini0647-queriedpositionprobability, def:physics:phyx-mini-0647:phyxminiproblems-problemphyxmini0647-satisfiesbornpositionphysics, def:physics:phyx-mini-0647:phyxminiproblems-problemphyxmini0647-matchesproblemstatementandfigure, def:physics:phyx-mini-0647:phyxminiproblems-problemphyxmini0647-matchesanswerchoice}
The assigned autoformalize agent should translate this physics problem into Lean declarations in the covered file, with theorem and lemma proof bodies written as `by sorry`.
\end{theorem}
\begin{proof}
This is an autoformalization task, not a proof task. Produce statements that can later be proved by the physics prover without weakening the physical model.
\end{proof}

% --- Archon declaration topology begin ---
\section{Declaration topology}
\begin{definition}[Quantum Particle Kind]
  \label{def:physics:phyx-mini-0647:phyxminiproblems-problemphyxmini0647-quantumparticlekind}
  \lean{PhyXMiniProblems.ProblemPhyXMini0647.QuantumParticleKind}
  The particle species stated in the problem
\end{definition}
\begin{proof}
  This is the declared model definition of Quantum Particle Kind.
\end{proof}
\begin{definition}[Horizontal Axis Quantity]
  \label{def:physics:phyx-mini-0647:phyxminiproblems-problemphyxmini0647-horizontalaxisquantity}
  \lean{PhyXMiniProblems.ProblemPhyXMini0647.HorizontalAxisQuantity}
  Physical quantity named on the horizontal axis
\end{definition}
\begin{proof}
  This is the declared model definition of Horizontal Axis Quantity.
\end{proof}
\begin{definition}[Position Axis Unit]
  \label{def:physics:phyx-mini-0647:phyxminiproblems-problemphyxmini0647-positionaxisunit}
  \lean{PhyXMiniProblems.ProblemPhyXMini0647.PositionAxisUnit}
  Unit printed on the horizontal position axis
\end{definition}
\begin{proof}
  This is the declared model definition of Position Axis Unit.
\end{proof}
\begin{definition}[Vertical Axis Quantity]
  \label{def:physics:phyx-mini-0647:phyxminiproblems-problemphyxmini0647-verticalaxisquantity}
  \lean{PhyXMiniProblems.ProblemPhyXMini0647.VerticalAxisQuantity}
  Physical quantity named by the vertical-axis label ψ(x)
\end{definition}
\begin{proof}
  This is the declared model definition of Vertical Axis Quantity.
\end{proof}
\begin{definition}[Position Tick]
  \label{def:physics:phyx-mini-0647:phyxminiproblems-problemphyxmini0647-positiontick}
  \lean{PhyXMiniProblems.ProblemPhyXMini0647.PositionTick}
  The four position ticks visible in the image
\end{definition}
\begin{proof}
  This is the declared model definition of Position Tick.
\end{proof}
\begin{definition}[Amplitude Tick]
  \label{def:physics:phyx-mini-0647:phyxminiproblems-problemphyxmini0647-amplitudetick}
  \lean{PhyXMiniProblems.ProblemPhyXMini0647.AmplitudeTick}
  The four amplitude ticks visible on the vertical axis
\end{definition}
\begin{proof}
  This is the declared model definition of Amplitude Tick.
\end{proof}
\begin{definition}[Wave Function Curve Shape]
  \label{def:physics:phyx-mini-0647:phyxminiproblems-problemphyxmini0647-wavefunctioncurveshape}
  \lean{PhyXMiniProblems.ProblemPhyXMini0647.WaveFunctionCurveShape}
  Qualitative shape of the red curve in the supplied image
\end{definition}
\begin{proof}
  This is the declared model definition of Wave Function Curve Shape.
\end{proof}
\begin{definition}[expected Position Tick Nanometers]
  \label{def:physics:phyx-mini-0647:phyxminiproblems-problemphyxmini0647-expectedpositionticknanometers}
  \lean{PhyXMiniProblems.ProblemPhyXMini0647.expectedPositionTickNanometers}
  \uses{def:physics:phyx-mini-0647:phyxminiproblems-problemphyxmini0647-positiontick}
  Nanometre coordinate readout printed at each horizontal tick
\end{definition}
\begin{proof}
  This is the declared model definition of expected Position Tick Nanometers.
\end{proof}
\begin{definition}[expected Amplitude Tick Multiplier]
  \label{def:physics:phyx-mini-0647:phyxminiproblems-problemphyxmini0647-expectedamplitudetickmultiplier}
  \lean{PhyXMiniProblems.ProblemPhyXMini0647.expectedAmplitudeTickMultiplier}
  \uses{def:physics:phyx-mini-0647:phyxminiproblems-problemphyxmini0647-amplitudetick}
  Each vertical tick as a dimensionless multiple of the positive scale c
\end{definition}
\begin{proof}
  This is the declared model definition of expected Amplitude Tick Multiplier.
\end{proof}
\begin{definition}[Piecewise Wave Function Figure]
  \label{def:physics:phyx-mini-0647:phyxminiproblems-problemphyxmini0647-piecewisewavefunctionfigure}
  \lean{PhyXMiniProblems.ProblemPhyXMini0647.PiecewiseWaveFunctionFigure}
  \uses{def:physics:phyx-mini-0647:phyxminiproblems-problemphyxmini0647-horizontalaxisquantity, def:physics:phyx-mini-0647:phyxminiproblems-problemphyxmini0647-positionaxisunit, def:physics:phyx-mini-0647:phyxminiproblems-problemphyxmini0647-verticalaxisquantity, def:physics:phyx-mini-0647:phyxminiproblems-problemphyxmini0647-positiontick, def:physics:phyx-mini-0647:phyxminiproblems-problemphyxmini0647-amplitudetick, def:physics:phyx-mini-0647:phyxminiproblems-problemphyxmini0647-wavefunctioncurveshape}
  The labeled numerical data in the wavefunction plot. amplitudeScalePerSqrtNanometer is the real readout denoted by c; it is not a scalar replacement for the electron state itself
\end{definition}
\begin{proof}
  This is the declared model definition of Piecewise Wave Function Figure.
\end{proof}
\begin{definition}[Electron Position Experiment]
  \label{def:physics:phyx-mini-0647:phyxminiproblems-problemphyxmini0647-electronpositionexperiment}
  \lean{PhyXMiniProblems.ProblemPhyXMini0647.ElectronPositionExperiment}
  \uses{def:physics:phyx-mini-0647:phyxminiproblems-problemphyxmini0647-quantumparticlekind, def:physics:phyx-mini-0647:phyxminiproblems-problemphyxmini0647-piecewisewavefunctionfigure}
  The electron-position experiment. The wavefunction is a complex-valued coordinate representative, while positionProbability is an independent normalized measure until Born's law is imposed below
\end{definition}
\begin{proof}
  This is the declared model definition of Electron Position Experiment.
\end{proof}
\begin{definition}[plotted Wave Amplitude Per Sqrt Nanometer]
  \label{def:physics:phyx-mini-0647:phyxminiproblems-problemphyxmini0647-plottedwaveamplitudepersqrtnanometer}
  \lean{PhyXMiniProblems.ProblemPhyXMini0647.plottedWaveAmplitudePerSqrtNanometer}
  \uses{def:physics:phyx-mini-0647:phyxminiproblems-problemphyxmini0647-piecewisewavefunctionfigure}
  The piecewise-constant amplitude shown in the primary image. Values at the four jump points use a fixed half-open convention; changing finitely many endpoint values does not change any Born probability
\end{definition}
\begin{proof}
  This is the declared model definition of plotted Wave Amplitude Per Sqrt Nanometer.
\end{proof}
\begin{definition}[queried Region Nanometers]
  \label{def:physics:phyx-mini-0647:phyxminiproblems-problemphyxmini0647-queriedregionnanometers}
  \lean{PhyXMiniProblems.ProblemPhyXMini0647.queriedRegionNanometers}
  \uses{def:physics:phyx-mini-0647:phyxminiproblems-problemphyxmini0647-piecewisewavefunctionfigure}
  The closed interval from the -1 tick to the 1 tick, in nanometre readouts
\end{definition}
\begin{proof}
  This is the declared model definition of queried Region Nanometers.
\end{proof}
\begin{definition}[queried Position Probability]
  \label{def:physics:phyx-mini-0647:phyxminiproblems-problemphyxmini0647-queriedpositionprobability}
  \lean{PhyXMiniProblems.ProblemPhyXMini0647.queriedPositionProbability}
  \uses{def:physics:phyx-mini-0647:phyxminiproblems-problemphyxmini0647-electronpositionexperiment, def:physics:phyx-mini-0647:phyxminiproblems-problemphyxmini0647-queriedregionnanometers}
  Dimensionless probability that the electron position lies in the queried region
\end{definition}
\begin{proof}
  This is the declared model definition of queried Position Probability.
\end{proof}
\begin{definition}[Satisfies Born Position Physics]
  \label{def:physics:phyx-mini-0647:phyxminiproblems-problemphyxmini0647-satisfiesbornpositionphysics}
  \lean{PhyXMiniProblems.ProblemPhyXMini0647.SatisfiesBornPositionPhysics}
  \uses{def:physics:phyx-mini-0647:phyxminiproblems-problemphyxmini0647-electronpositionexperiment}
  The square-integrability requirement and Born's position-probability law. Lebesgue volume is taken on the numerical nanometre coordinate, so the squared amplitude is embedded as a nonnegative density with ENNReal.ofReal
\end{definition}
\begin{proof}
  This is the declared model definition of Satisfies Born Position Physics.
\end{proof}
\begin{definition}[Matches Problem Statement And Figure]
  \label{def:physics:phyx-mini-0647:phyxminiproblems-problemphyxmini0647-matchesproblemstatementandfigure}
  \lean{PhyXMiniProblems.ProblemPhyXMini0647.MatchesProblemStatementAndFigure}
  \uses{def:physics:phyx-mini-0647:phyxminiproblems-problemphyxmini0647-expectedpositionticknanometers, def:physics:phyx-mini-0647:phyxminiproblems-problemphyxmini0647-expectedamplitudetickmultiplier, def:physics:phyx-mini-0647:phyxminiproblems-problemphyxmini0647-electronpositionexperiment, def:physics:phyx-mini-0647:phyxminiproblems-problemphyxmini0647-plottedwaveamplitudepersqrtnanometer}
  Scenario facts and exact readouts from 647.png. These assumptions contain the plotted state but no probability for the central interval and no answer-choice value
\end{definition}
\begin{proof}
  This is the declared model definition of Matches Problem Statement And Figure.
\end{proof}
\begin{definition}[Answer Choice]
  \label{def:physics:phyx-mini-0647:phyxminiproblems-problemphyxmini0647-answerchoice}
  \lean{PhyXMiniProblems.ProblemPhyXMini0647.AnswerChoice}
  The four answer labels printed in the problem
\end{definition}
\begin{proof}
  This is the declared model definition of Answer Choice.
\end{proof}
\begin{definition}[displayed Probability]
  \label{def:physics:phyx-mini-0647:phyxminiproblems-problemphyxmini0647-answerchoice-displayedprobability}
  \lean{PhyXMiniProblems.ProblemPhyXMini0647.AnswerChoice.displayedProbability}
  \uses{def:physics:phyx-mini-0647:phyxminiproblems-problemphyxmini0647-answerchoice}
  Dimensionless probability displayed beside each answer label
\end{definition}
\begin{proof}
  This is the declared model definition of displayed Probability.
\end{proof}
\begin{definition}[Matches Answer Choice]
  \label{def:physics:phyx-mini-0647:phyxminiproblems-problemphyxmini0647-matchesanswerchoice}
  \lean{PhyXMiniProblems.ProblemPhyXMini0647.MatchesAnswerChoice}
  \uses{def:physics:phyx-mini-0647:phyxminiproblems-problemphyxmini0647-electronpositionexperiment, def:physics:phyx-mini-0647:phyxminiproblems-problemphyxmini0647-queriedpositionprobability, def:physics:phyx-mini-0647:phyxminiproblems-problemphyxmini0647-answerchoice}
  A displayed choice matches the modeled probability exactly
\end{definition}
\begin{proof}
  This is the declared model definition of Matches Answer Choice.
\end{proof}
\begin{lemma}[central interval probability]
  \label{lem:physics:phyx-mini-0647:phyxminiproblems-problemphyxmini0647-central-interval-probability}
  \lean{PhyXMiniProblems.ProblemPhyXMini0647.central_interval_probability}
  \uses{def:physics:phyx-mini-0647:phyxminiproblems-problemphyxmini0647-electronpositionexperiment, def:physics:phyx-mini-0647:phyxminiproblems-problemphyxmini0647-queriedpositionprobability, def:physics:phyx-mini-0647:phyxminiproblems-problemphyxmini0647-satisfiesbornpositionphysics, def:physics:phyx-mini-0647:phyxminiproblems-problemphyxmini0647-matchesproblemstatementandfigure}
  The normalized piecewise wavefunction assigns probability 4/5 to the interval from -1 nm through 1 nm
\end{lemma}
\begin{proof}
  The conclusion follows from the governing relations and figure readouts named in the statement of central interval probability.
\end{proof}
% --- Archon declaration topology end ---
% --- Archon physics formalization source end ---
